% 1 章：はじめに
% 数値の出典：memo/09_results_synthetic.md（S1〜S4）、memo/10_results_empirical.md（R1〜R4）、docs/tables/*.tex
\section{はじめに}\label{sec:intro}

\subsection{背景：ボラティリティ推定の実務上の位置付け}\label{sec:intro-background}

金融機関の業務システムにおいて、資産価格の変動の大きさ、すなわちボラティリティは、最も基本的で、かつ最も頻繁に使われる入力の一つである。
ボラティリティの推定値を入力として使う業務は多い。主なものを挙げる。

\begin{itemize}
  \item \textbf{市場リスク管理}：
    バリュー・アット・リスク（VaR）や期待ショートフォール（ES）は、保有ポートフォリオの損益分布の裾を表す指標であり、その計算はリターンの条件付き分布、とりわけその散らばりの推定から始まる。
    バーゼル銀行監督委員会による市場リスクの最低所要自己資本の枠組み \parencite{bcbs2019} では、内部モデル方式で使うリスク指標が 99\% VaR から 97.5\% ES に改められ、損益分布の裾をより重視する方向に進んでいる。
    裾の推定は、ボラティリティの推定の質に直接左右される。
  \item \textbf{デリバティブの評価とヘッジ}：
    オプションの価格やヘッジ比率（デルタなど）は将来のボラティリティに依存する。
    市場で観測されるインプライド・ボラティリティと、過去のリターンから推定したボラティリティを比べることは、値付けの妥当性の確認やトレーディングの判断材料になる。
  \item \textbf{ポートフォリオの構築と運用}：
    ボラティリティ・ターゲティング（ボラティリティが上がったらポジションを落とす）、リスクパリティ（各資産のリスク寄与を揃える）、個別銘柄のポジションサイズの決定などでは、ボラティリティの推定値がそのまま売買の量を決める。
  \item \textbf{証拠金・担保・与信の管理}：
    清算機関や取引相手に差し入れる証拠金、担保の掛け目、市場性与信の枠などは、価格変動の大きさに応じて決まる。
  \item \textbf{モニタリングとアラート}：
    相場の急変を検知し、リスク管理部門やフロントに通知する仕組みでも、「いまのボラティリティが平常時と比べてどの程度高いか」が判断の基準になる。
\end{itemize}

ボラティリティが時間を通じて一定ではないことは、古くからよく知られている。
大きな価格変動の後には大きな変動が続き、穏やかな時期は穏やかな時期が続く（ボラティリティ・クラスタリング）。
また、日次リターンの分布は正規分布よりも裾が厚い。
本稿で使う実データ（日経平均株価と TOPIX に連動する ETF の 2024 年 9 月から 2026 年 8 月までの日次リターン、\cref{sec:empirical}）でも、
年率換算のボラティリティ（事後平均）は期間の中央値で 16〜21\% 程度であるのに対し、
2025 年 4 月の米国の関税政策をきっかけとした急落・急反発の局面では、数営業日のうちに 57〜61\% まで跳ね上がった。
こうした急変に素早く追随しつつ、平常時には過剰に反応しないことが、ボラティリティの推定に求められる。

システムを作る立場から見ると、ボラティリティの推定は一度だけ行う分析ではなく、
\textbf{多数の銘柄について、毎営業日、決められた時刻までに}繰り返し実行されるバッチ処理である。
例えば、夜間バッチで数百〜数千の銘柄のボラティリティを推定し、その結果を翌朝の寄り付き前までにリスク計算やレポートに渡す、といった形になる。
推定結果は下流のリスク計算・規制報告・顧客向けの帳票などに流れるため、推定処理には次のような要件が課される。

\begin{itemize}
  \item \textbf{処理時間が予測できること}：バッチの締め切りに間に合うかどうかを事前に見積もれること。銘柄数やデータ長が増えたときの伸び方が分かっていること。
  \item \textbf{失敗しにくく、失敗を検知できること}：推定がうまくいかなかったときに、それを自動で検知し、再実行やフォールバックができること。
  \item \textbf{結果が再現できること}：同じ入力と設定からは同じ結果が得られ、後から監査や検証ができること。
  \item \textbf{説明できること}：推定値がなぜその値になったのかを、リスク管理部門・監査・監督当局に説明できること。とくに、リスクを小さく見積もる方向の偏りは強く問題視される。
\end{itemize}

これらの要件は、どのモデルを使うかだけでなく、\textbf{そのモデルをどのような計算方法で推定するか}の選択を大きく左右する。
これが本稿の問題意識の出発点である。

\subsection{GARCH 型モデルと確率的ボラティリティモデル}\label{sec:intro-models}

時間変動するボラティリティを表す統計モデルには、大きく二つの系統がある。

一つは GARCH 型のモデルである。代表的な GARCH(1,1) モデル \parencite{bollerslev1986} では、日次リターン $y_t$ の条件付き分散 $\sigma_t^2$ を
\[
  \sigma_t^2 = \omega + \alpha\, y_{t-1}^2 + \beta\, \sigma_{t-1}^2
\]
と、前日のリターンの二乗と前日の条件付き分散の線形結合で表す。
条件付き分散は過去の観測から\textbf{決定論的に}決まるので、尤度は各日の条件付き密度の積として逐次的に計算でき、最尤法で容易に推定できる。
推定は 1 銘柄あたりミリ秒単位で終わり、解釈もしやすいため、GARCH 型モデルは実務の標準になっている。
本稿でも、実データでの予測性能の比較対象として GARCH(1,1) を使う（\cref{sec:empirical}）。

もう一つが、本稿で扱う\textbf{確率的ボラティリティ（stochastic volatility, SV）モデル} \parencite{taylor1986} である。
基本的な SV モデルは、対数ボラティリティ $h_t$ を観測されない状態変数として
\[
  y_t = \exp(h_t/2)\, \varepsilon_t, \qquad
  h_{t+1} = \mu + \phi\,(h_t - \mu) + \sigma\, \eta_t, \qquad
  \varepsilon_t,\ \eta_t \sim \Normal(0, 1)
\]
と表す（詳しくは\cref{sec:model}）。
$\mu$ は対数ボラティリティの長期的な水準、$\phi$ はその持続性、$\sigma$ はボラティリティ自体の変動の大きさを表す。
GARCH との本質的な違いは、ボラティリティが観測 $y_t$ とは別の固有のショック $\eta_t$ を持つ点にある。
ボラティリティは過去のリターンから決まるのではなく、それ自体が確率的に動く隠れた過程であり、リターンはその過程の「ノイズを含んだ観測」とみなされる。

SV モデルには、実務の観点からもいくつかの利点がある。

\begin{itemize}
  \item \textbf{解釈が自然である}：
    市場に到着する情報の量が日々確率的に変動し、それに応じて価格変動の大きさが変わる、という見方（混合分布仮説、\cite{clark1973}）を素直にモデル化したものと解釈できる。
    また、オプション評価で使われる連続時間の確率ボラティリティモデルを離散時間で表したものとも見なせる。
  \item \textbf{ボラティリティ自体の不確実性を表せる}：
    ベイズ推定では、パラメータ $(\mu,\phi,\sigma)$ だけでなく、各日の対数ボラティリティ $h_t$ についても事後分布が得られる。
    例えば「今日のボラティリティは 25\% である」という点推定だけでなく、「95\% の確からしさで 20〜32\% の範囲にある」という区間として示せることは、リスク管理部門や運用担当者との対話で有用である。
  \item \textbf{拡張が自然である}：
    株価が下落するとボラティリティが上がる傾向（レバレッジ効果）は $\varepsilon_t$ と $\eta_t$ の相関として \parencite{omori2007}、
    ボラティリティの変動だけでは説明しきれない裾の厚さは $\varepsilon_t$ を $t$ 分布にすることで、それぞれ同じ枠組みの中で表せる。
    本稿ではこの二つの拡張も扱う。
  \item \textbf{当てはまりと予測の良さ}：
    SV モデルは ARCH 型モデルと比べて遜色のない、あるいはより良い当てはまりを示すことが報告されている \parencite{kim1998}。
    本稿の実データ分析でも、1 期先の予測分布の良さを測る対数予測スコアで、SV モデルが GARCH(1,1) を上回った（差は小さく、統計的に有意だったのは 2 銘柄のうち 1 銘柄。\cref{sec:empirical}）。
\end{itemize}

一方で、SV モデルには推定上の大きな難しさがある。
パラメータ $\theta = (\mu,\phi,\sigma)$ の尤度は、観測されない $h_1,\dots,h_T$ をすべて積分消去した
\[
  p(y_{1:T} \mid \theta) = \int \prod_{t=1}^{T} p(y_t \mid h_t)\; p(h_{1:T} \mid \theta)\; dh_{1:T}
\]
で与えられる。これは $T$ 次元（日次データ 2 年分なら約 500 次元、10 年分なら約 2,500 次元）の積分であり、閉じた形では計算できない。
GARCH 型モデルでは条件付き分散が観測から決まるためこの積分が現れず、ここが両者の計算コストの違いの根源である。
SV モデルの推定では、この高次元の積分をどう扱うかが常に中心的な問題になる。

\subsection{従来の推定法とシステム運用上の課題}\label{sec:intro-mcmc}

SV モデルの推定法は、これまで数多く提案されてきた。
初期の代表的な方法は、$\log y_t^2 = h_t + \log \varepsilon_t^2$ という変形により観測方程式を線形にし、$\log \varepsilon_t^2$ の分布を正規分布で近似したうえでカルマンフィルタを使う疑似最尤法（QML）である \parencite{harvey1994}。
計算は速いが、$\log \varepsilon_t^2$ の分布（裾が非対称な $\log \chi^2_1$ 分布）を正規分布で置き換えるため、推定の効率が落ちる。

その後、\textbf{マルコフ連鎖モンテカルロ法（MCMC）}によるベイズ推定が SV モデルの標準的な推定法になった。
\textcite{jacquier1994} は、$h_t$ を一つずつ更新する MCMC で、パラメータと潜在ボラティリティの事後分布を同時に求める方法を示した。
\textcite{kim1998} は、$\log \varepsilon_t^2$ の分布を正規分布の混合で精度よく近似することで、$h_{1:T}$ 全体をカルマンフィルタ・スムーザーを使って一度にサンプリングする効率的な方法を提案した。
\textcite{omori2007} はこれをレバレッジ効果のあるモデルに拡張し、
\textcite{chan2009} は、状態変数の事後分布の精度行列（共分散行列の逆行列）が帯行列になることを使って、状態を $O(T)$ の計算量で一括してサンプリングする方法を示した。
\textcite{kastner2014} は、中心化・非中心化という二つのパラメータ化を交互に使う ASIS という手法で、とくに $\sigma$ の混合（サンプルの自己相関の小ささ）を大きく改善している。
日本語の解説としては \textcite{watanabe2000} がある。

近年は、モデルを確率的プログラミング言語で書けば勾配情報を使って自動的にサンプリングできる、汎用の MCMC 手法も普及した。
その代表が、ハミルトニアン・モンテカルロ法の歩幅や軌道の長さを自動で調整する No-U-Turn Sampler（NUTS）\parencite{hoffman2014} であり、Stan や NumPyro \parencite{phan2019} などのライブラリで使える。
本稿では、NumPyro の NUTS を十分に長く回した結果を、事後分布の「正解」として扱う。

MCMC は、反復を十分に長くとれば事後分布を正確に表すサンプルが得られる（漸近的に正確な）方法であり、研究の場では SV 推定の標準である。
しかし、前節で述べたシステムの要件に照らすと、次のような課題がある。

\begin{enumerate}[label=(\arabic*)]
  \item \textbf{計算時間が長く、しかも事前に読みにくい}：
    MCMC のサンプルは互いに相関しているため、事後分布を十分な精度で表すには、得たいサンプル数の何倍もの反復が必要になる。
    その倍率はモデル・データ・パラメータ化によって大きく変わり、事前には分からない。
    本稿の計算時間の実験（\cref{sec:synthetic}）では、系列長 $T=10{,}000$ の SV モデルに NUTS を（ウォームアップ・サンプリング各 1,000 反復、4 チェーンで）適用すると 548 秒かかり、
    しかもパラメータの有効サンプルサイズ（ESS：自己相関を考慮した実質的なサンプル数）は最小で 14 しかなかった。
  \item \textbf{収束の判定と、判定に落ちたときの扱いが難しい}：
    MCMC の結果を使ってよいかは、複数チェーン間の一致度 $\hat R$、ESS、発散の有無などで判定する。
    判定自体は自動化できるが、基準を満たさなかったときに反復を延長するのか、パラメータ化を変えて再実行するのか、といった運用の設計が必要になる。
    本稿の実験でも、標準的な設定（ウォームアップ・サンプリング各 2,000 反復、4 チェーン）では、
    合成データの 200 回の推定のうち 156 回、実データの 32 銘柄のうち 11 銘柄で、パラメータの $\hat R$ が目安の 1.05 を超えた（\cref{sec:synthetic,sec:system}）。
    SV モデルの MCMC で $\sigma$ の混合が悪いことはよく知られた問題であり \parencite{kastner2014}、汎用の NUTS でもそれは解消されない。
  \item \textbf{銘柄数・再推定の頻度とともに計算量が積み上がる}：
    実データの 32 銘柄（各 485 日）を NUTS で 1 銘柄ずつ推定すると、合計で 293 秒かかった（\cref{sec:system}）。
    実務では、これが銘柄数（数百〜数千）、モデルの種類、毎営業日の再推定の分だけ掛け算で増える。
    前日の結果を初期値に使うことはできるが、MCMC ではウォームアップ（歩幅などの適応）とサンプリングそのものは毎回必要であり、計算量はあまり減らない。
\end{enumerate}

\subsection{変分推論という選択肢とその落とし穴}\label{sec:intro-vi}

こうした背景から、本稿では\textbf{変分推論（variational inference, VI）}に注目する \parencite{blei2017}。
VI は、事後分布 $p(\theta, h_{1:T} \mid y_{1:T})$ を、計算しやすい分布の族 $\mathcal{Q}$（例えば多変量正規分布の族）の中で最も近い分布 $q^\ast$ で近似する方法である。
近さはカルバック・ライブラー情報量 $\KL(q \,\|\, p)$ で測り、それを最小化する代わりに、等価な目的関数である証拠下界（ELBO）を最大化する（\cref{sec:inference}）。
つまり VI は、\textbf{積分（サンプリング）の問題を最適化の問題に置き換える}。
再パラメータ化による勾配の推定 \parencite{kingma2014} と自動微分を組み合わせると、モデルの対数密度を書くだけで確率的勾配法により最適化できる。
これを、制約のあるパラメータの変換まで含めて自動化したものが自動微分変分推論（ADVI）\parencite{kucukelbir2017} である。

システムの観点から見ると、VI には次のような利点がある。

\begin{itemize}
  \item \textbf{速い}：数千ステップの勾配法で済み、MCMC より桁違いに短い時間で事後分布の近似が得られる。
  \item \textbf{処理時間が予測しやすい}：計算量は「ステップ数 × 1 ステップのコスト」で決まり、ステップ数は設定で固定できる。収束の判定も、目的関数（ELBO）の推移という一つの数値列を見ればよい。
  \item \textbf{前日の結果を再利用できる}：最適化なので、前日の最適解から再開（warm start）すれば、少ないステップで今日の解に到達できる。
  \item \textbf{並列化・一括処理と相性が良い}：銘柄ごとの最適化は互いに独立であり、そのまま複数の計算資源に振り分けられる。
  \item \textbf{結果を固定長の数値として保存できる}：推定結果は変分分布のパラメータ（数千個の数値）として保存でき、そこから必要なときに必要な数だけサンプルを生成できる。乱数のシードを固定すれば再現も容易である。
\end{itemize}

ただし、VI には重要な落とし穴がある。
第一に、VI は近似であり、得られる分布の精度の上限は、あらかじめ選んだ族 $\mathcal{Q}$ の表現力で決まる。反復を増やしても、族の外にある真の事後分布には近づかない。
第二に、VI が最小化する $\KL(q \,\|\, p)$ は非対称な量であり、真の事後分布 $p$ の確率が小さいところに $q$ が確率を置くことを強く罰する。
その結果、$q$ は $p$ の裾を切り捨てる傾向があり、\textbf{事後分布の分散（不確実性）を過小評価しやすい} \parencite{blei2017}。
リスク管理の文脈では、これは「リスクを実際より小さく見せる」方向の偏りであり、前節で述べた説明責任の観点から最も避けたい性質である。

とりわけ問題になるのが、最も単純で広く使われる\textbf{平均場近似}（すべての変数を互いに独立とみなす族）である。
SV モデルでは、対数ボラティリティ $h_t$ が $\phi \approx 0.95$〜$0.99$ という強い持続性を持ち、事後分布でも隣り合う日の $h_t$ どうしが強く相関する。
平均場近似はこの相関をまったく表せない。
本稿の実験では、平均場近似による VI は、真値が 0.97〜0.99 のシナリオで持続性 $\phi$ の推定値を 0.03〜0.3 程度まで崩してしまい、
「分散を過小評価する」という教科書的な問題以前に、\textbf{パラメータの点推定そのものが大きく偏る}ことが分かった（\cref{sec:synthetic}）。
一方で、表現力を上げるために変数間の相関をすべて許すフルランクの正規分布（変分パラメータは約 50 万個）を使っても、
最適化が難しくなり、真値 $\phi = 0.99$、$T = 1{,}000$ のデータに 230 秒かけても $\phi = 0.24$ にとどまった。
\textbf{パラメータを増やせばよいわけではなく、モデルの構造に合った族を選ぶことが重要}なのである。

したがって、SV モデルに VI を使えるかどうかは、
「SV の事後分布の構造を正しく捉えつつ、VI の速さと運用上の利点を保てる変分族を設計できるか」にかかっている。

\subsection{関連研究と本稿の位置付け}\label{sec:intro-related}

時系列の状態空間モデルに対する VI では、状態変数の時間方向の依存を変分族に組み込む研究が進んできた。
\textcite{archer2015} は、時間方向にブロック三重対角の精度行列を持つ正規分布を変分族とし、カルマンフィルタ・スムーザーと同じ直観で、非線形な状態空間モデルにも使える汎用の VI を構成した。
\textcite{tan2018} は、精度行列のコレスキー因子にモデルの条件付き独立性に応じた疎な構造を課した正規分布による変分近似を提案し、状態空間モデルに適用している。
これらは、MCMC の側で \textcite{chan2009} が使った「帯行列の精度行列による $O(T)$ の計算」と同じ構造を、変分族の側で利用するものである。

さらに、パラメータ（大域変数）と状態（局所変数）の事後的な依存を変分族に取り入れる研究もある。
\textcite{tan2020} は、階層モデルでは分散などの大域パラメータが局所潜在変数の事後分布の尺度を決めることに着目し、局所変数の変分分布を大域パラメータに条件付けた族を提案して、一般化線形混合モデルや状態空間モデルに適用した。
\textcite{loaizamaya2022} は、パラメータには簡潔なパラメトリック近似を使い、潜在変数には（近似ではなく）正確な条件付き事後分布を使う族を提案し、その密度を計算できなくても、条件付き事後分布からサンプリングさえできれば勾配が計算できることを示した。
\textcite{frazier2023} は、状態空間モデルでよく使われる VI の方法を比較し、状態の事後分布を近似しない方法の方が、状態まで近似する方法よりもパラメータの推定が正確であること、
ただし推論の精度の差は短い検証期間では予測精度の小さな差にしかならないことが多い一方、長い検証期間では意味のある差になりうることを報告している。
機械学習の分野では、\textcite{johnson2016} が、状態の事後分布を「事前分布 × ニューラルネットワークが出力する共役な因子（擬似観測）」の形で表す構造化 VI を提案している。

SV モデルそのものに VI を適用した研究としては、
多変量の因子 SV モデルに対する高速な変分ベイズ法と、新しい観測が来たときの逐次的な変分更新を提案した \textcite{gunawan2021}、
多数の変数を持つベイズ VAR の SV 部分について、対数ボラティリティの同時事後分布を大域的に近似する方法を提案した \textcite{chan2022}、
状態の同時事後分布を動的因子モデルで近似する正規分布の変分近似を、Wishart 型の多変量 SV モデルに適用した \textcite{quiroz2023} などがある。

本稿で提案として扱う\textbf{条件付き構造化変分族}（\cref{sec:family}）は、
状態 $h_{1:T}$ の変分分布をパラメータ $\theta$ に条件付け、その形を「AR(1) の事前分布 × 正規分布の擬似観測」とするものである。
これは、上に挙げた考え方、とくに \textcite{tan2020} の条件付きの構造と \textcite{johnson2016} の擬似観測による表現を、一変量の SV モデルに素直に当てはめたものであり、
手法として新しいことを主張するものではない。
本稿の主眼はむしろ、次の三点にある。

\begin{itemize}
  \item 平均場近似から条件付き構造化まで、\textbf{変分族の設計の違いが結果をどれほど変えるか}を、同じ実装・同じ条件のもとで、真値（合成データ）と十分に長い MCMC を基準に定量的に示すこと。
  \item 日本の株式市場の実データ（ETF と個別株）で、事後分布の一致だけでなく、\textbf{VaR の予測性能まで}評価すること。
  \item システムに載せるときの具体的な論点（日次の再推定、多銘柄の並列化、結果の監視、MCMC との使い分け）を、\textbf{実測に基づいて}整理すること。
\end{itemize}

\subsection{本稿の目的と貢献}\label{sec:intro-contrib}

本稿の目的は、SV モデルのベイズ推定を VI で行うことが金融機関向けシステムの実務に使えるかを、\textbf{精度・計算時間・運用性}の三つの観点から評価し、
使えるとすればどのように使うべきかを示すことである。
本稿の貢献は次の四つにまとめられる。

\begin{enumerate}[label=(\roman*)]
  \item \textbf{SV モデルに対する変分族の体系的な比較}（\cref{sec:family,sec:synthetic}）。
    基本 SV・レバレッジ SV・$t$ 誤差 SV の三つのモデルについて、
    平均場、低ランク、フルランク、構造化（パラメータと状態を独立とし、状態は三重対角の精度行列を持つ正規分布）、条件付き構造化の各変分族を、
    真値が分かっている合成データ（5 つのシナリオ × 2 つの系列長 × 各 20 本）と、十分に長く回した NUTS を基準に比較した。
    その結果、平均場・低ランク・フルランクの族はいずれも持続性 $\phi$ の推定が崩れること、
    構造化の族は状態 $h_t$ の推定は良好だが $\sigma$ の不確実性を 1/5〜1/10 に過小評価し、$\sigma$ の 95\% 区間が真値をほとんど含まないことを示した。
  \item \textbf{条件付き構造化変分族の構成と評価}（\cref{sec:family,sec:synthetic}）。
    $q(\theta, h) = q(\theta)\, q(h \mid \theta)$、
    $q(h \mid \theta) \propto p(h \mid \theta) \prod_{t} \Normal(\tilde y_t \mid h_t, \lambda_t^{-1})$
    という族を、SV モデル向けに具体的に構成した（$\tilde y_t$ と $\lambda_t$ が変分パラメータ）。
    $\theta$ を固定すると $q(h \mid \theta)$ の精度行列は三重対角になり、サンプリングも密度の計算も $O(T)$ で済む。
    また、この族は中心化・非中心化というパラメータ化の選び方に依存しない。レバレッジ SV への拡張も示した。
    持続性の高いシナリオ（$\phi = 0.99$、$T = 1{,}000$）で、事後平均は NUTS とほぼ一致し（$\phi$：0.983 対 0.983、$\sigma$：0.123 対 0.126）、
    $h_t$ の事後標準偏差も NUTS の 0.98 倍と、不確実性の大きさまで再現できた。
    計算時間は $T = 10{,}000$ で約 11 秒であり、NUTS（548 秒）の約 50 倍速い。
  \item \textbf{日本の株式市場の実データでの評価}（\cref{sec:empirical}）。
    J-Quants API の無料プランで取得した、日経平均株価・TOPIX に連動する ETF（1321・1306）と TOPIX Core30 構成銘柄の日次リターン（485 営業日）を使った。
    全期間の推定では、提案手法の事後分布は NUTS とほぼ一致した（例：1321 の $\phi$ は 0.898 対 0.889）。
    250 日の窓を 1 日ずつずらしながら毎日推定し直す VaR のバックテスト（検証 235 日）では、
    提案手法による SV モデルの 5\% VaR の超過率は 6.4〜6.8\% で、Kupiec の検定・Christoffersen の検定のいずれでも棄却されず、
    対数予測スコアは GARCH(1,1)（正規誤差・$t$ 誤差のいずれ）よりも良かった（ただし差は 1 日あたり 0.03 nats 以下と小さく、Diebold–Mariano 検定で有意だったのは 1306 だけである）。
    一方、平均場近似では 1321 の 5\% VaR の超過率が 8.1\% となり、有意水準 5\% の検定で棄却された。
  \item \textbf{システム実装の観点からの知見}（\cref{sec:system}）。
    前日の変分パラメータを 1 日ずらして初期値にする warm start により、毎日の再推定は 1 銘柄 1 日あたり 0.10 秒（毎回ゼロから推定する場合の約 1/5）になり、VaR への影響は中央値で 0.02〜0.03\% ポイントと無視できる大きさだった。
    32 銘柄の一括推定は、NUTS の逐次実行で 293 秒、提案手法の逐次実行で 50 秒、8 つの CPU デバイスへの並列実行で 15 秒（NUTS の約 19 倍速い）だった。
    また、銘柄の方向に計算をベクトル化（\code{jax.vmap}）すると 68 秒と、逐次実行よりかえって遅くなることを確認した。
    「ベクトル化すれば速くなる」とは限らず、計算の形とハードウェアに合った並列化を選ぶ必要がある。
    これらを踏まえ、日次の運用は warm start の VI で行い、定期的に MCMC で結果を突き合わせる\textbf{ハイブリッド運用}と、その監視の方法を提案する。
\end{enumerate}

同時に、本稿の結果には限界もある。
提案手法は、レバレッジ SV の相関パラメータ $\rho$ を 0 の側に約 25\% 縮めて推定する。
実データでは、$\sigma$ を NUTS より 8\% ほど小さく、その事後標準偏差も NUTS の約 0.85 倍と、不確実性をやや小さく見積もる傾向が見られた。
VI の当てはまりを自動で診断する指標として提案されている PSIS の $\hat k$ は、手法の良し悪しの順位とはおおむね整合したものの、本稿の問題では提案手法を含むすべての手法で目安の 0.7 を超えた（提案手法でも 0.85）。そのため、「この値を下回れば近似を信頼してよい」という単独の判定基準としては使いにくい。
また、無料プランで取得できるデータは約 2 年分と短く、VaR の検証期間（235 日）では 1\% VaR の検定の検出力が低い。
計算時間はすべて単一の CPU マシン（12 コア）で測ったものであり、GPU での評価は行っていない。
これらの点は\cref{sec:synthetic,sec:empirical,sec:system,sec:conclusion}で詳しく述べる。

\subsection{本稿の構成}\label{sec:intro-outline}

本稿の構成は次のとおりである。

\cref{sec:model}では、基本 SV モデルと、その二つの拡張（レバレッジ効果、$t$ 分布の誤差）を定義し、事前分布を示す。
あわせて、尤度が解析的に書けない理由と、推定上の論点になる中心化・非中心化のパラメータ化を説明する。

\cref{sec:inference}では、ベイズ推定の枠組みと、本稿で「正解」として使う NUTS の概要を述べたうえで、
VI の基礎（KL 情報量と ELBO の関係、再パラメータ化による勾配の推定、ADVI）を、読者が実装の見通しを持てる程度に詳しく説明する。

\cref{sec:family}は本稿の中心であり、SV モデルの事後分布の構造（AR(1) の事前分布の精度行列が三重対角であること）を確認したうえで、
平均場から条件付き構造化までの変分族を定義し、それぞれのサンプリング・エントロピーの計算方法と計算量を比較する。

\cref{sec:synthetic}では、合成データを使って、パラメータの復元精度、NUTS との事後分布の一致度、系列長に対する計算時間の伸び方、設定への感度を評価する。

\cref{sec:empirical}では、J-Quants API で取得した ETF のデータを使って、全期間のボラティリティの推定と、ローリング推定による VaR のバックテストを行い、GARCH(1,1) と比較する。

\cref{sec:system}では、システムに組み込むことを想定して、日次の再推定（warm start）、多銘柄の並列処理、収束の判定と監視、再現性、MCMC との使い分けを、実測の結果に基づいて論じる。

\cref{sec:conclusion}では、結論と今後の課題（多変量 SV、実現ボラティリティの利用、より表現力の高い変分族など）を述べる。

付録には、ELBO の勾配と条件付き構造化変分族の導出（\cref{app:derivation}）、実装の抜粋（\cref{app:code}）、実験設定と感度分析（\cref{app:settings}）、データの取得と前処理（\cref{app:data}）をまとめた。

なお、本稿ではリターンを日次の対数リターンの 100 倍（\% 単位）で表し、標本平均を除いたものを $y_t$ とする。
年率換算のボラティリティは $\sqrt{252}\, \exp(h_t/2)$（\%）で表す。
